\scanpage{61}\mspage{60}
de la façon suivante : ce sont les $A\in M$ tels que $A=\bigcup_i A_{x_i}$, où \deleted{les} $\{x_i\}$ \deleted{\unsure{sont|soit|font}}\inserted{est} un ens fini\unsure{.} d'éléments de $M$, mutuellement compatibles (\deleted{\unsure{pour}}\inserted{au sens de} $\unsure{R|\deleted{R}}$.).

Donc une structure $(\mathcal F,\leq)$ satisfaisant $\deleted{\unsure{C_1}}C1,\deleted{\unsure{C_2}}C2,\deleted{\unsure{C_3}}C3$ + figures \deleted{\unsure{toutes|double|de de}} de type fini, équivaut à $(M,\leq,R)$, où \inserted{$R$ \uline{rel. sym. refl.}} \deleted{$(M,\unsure{\leq})$}\inserted{où $\leq$ et $R$} satisfont \deleted{deux} conditions.

a) \inserted{(concernant $\leq$)} Pour tte partie fermée \uline{majorée} $A$\deleted{\unsure{,}} de $M$, $A$ est de type fini i.e de la forme $\bigcup_{\deleted{\unsure{x_i|x}}i}\deleted{A_{\unsure{x_i|x}}}M_{\leq\unsure{x_i|x}}$, où les $x_i\in A$ \deleted{\unsure{fo}} sont un ens fini [ pour \unsure{ceci}, il suffit que les $M_{\leq x}$ soient finis ]

% slanted
\margin{donc \unsure{tte|toute|les} \unsure{partie|parties} \unsure{fermée|fermées} d'une partie fermée de type fini est de type fini}

b) \unsure{Compat.|(Compat.} de $\leq$ \uline{et} $R$) Si $(x,y)\in R$, et $x'\leq x$ $y'\leq y$, alors $\deleted{R}(x',y')\in R$.

% slanted
\margin{C'est \unsure{la|le} \unsure{propriété} Cont 2 de \unsure{pg|page|p°} 33}

(les ``figures'' de $M$ sont alors les parties \deleted{\unsure{fermées|fermée}}\inserted{$A$} de $M$ qui ont les propriétés suivantes :
\[
\begin{aligned}
\text{a) } & A\text{ est fermé}\\
\text{b) } & A\text{ est de type fini}\\
\text{c) } & \deleted{\unsure{A}}\unsure{\forall|\deleted{\forall}}\,x,y\in A,\ \text{on a}\quad (x,y)\in R.
\end{aligned}
\]
